\documentclass[12pt]{article}
%^ Start of document
\usepackage{times}  %Makes text TimesNewRoman
\usepackage{setspace} %Allows you to set single or double space 
\usepackage{biblatex} %Imports biblatex package
\usepackage{graphicx} % Required for inserting images
\usepackage{authblk}
\usepackage[colorlinks=true, urlcolor=blue, linkcolor=blue]{hyperref}
\usepackage{subfig}
\usepackage{float}
\usepackage{tabularx}
\usepackage{multirow}
\usepackage[paperheight=11in,
   paperwidth=8.5in,
   top=20mm,
   bottom=20mm,
   left=40mm,
   right=40mm]{geometry}
\usepackage{moresize}
\usepackage{tikz}
\usepackage{xcolor}
\usepackage{physics}
\usepackage{amssymb}
\usepackage{mathrsfs}
\usepackage{amsmath}
\usepackage[document]{ragged2e}
\usepackage{comment}
\usepackage{xcolor}
\addbibresource{mm.bib}
\begin{document}
\singlespacing  %Sets single spacing for whole document
\title{Suppression of radiation emissions from relativistic plasma by utilizing magnetic perturbations}

\author[1]{Zachary Tullier}
\affil[1]{Student\\Physics Department\\ University of Houston - Clear Lake \\Houston, Texas\\U.S}
\date{}
\maketitle

\section{Introduction}
    Atomic magnetic moments arise from the motion and spin of electrons. In an atom, each electron contributes an orbital magnetic moment (from its orbital angular momentum) and a spin magnetic moment (from its intrinsic spin). Quantum mechanically, the orbital magnetic moment operator is $\hat{\mu}_L = -\frac{e}{2m}\hat{L}$, so its $z$-component has eigenvalues $- \frac{e\hbar}{2m} m_l \equiv -\mu_B m_l$ (defining $\mu_B = \frac{e\hbar}{2m}$ as the Bohr magneton) \cite{PerturbationTheory}. Similarly, the electron’s spin contributes a moment $\hat{\mu}_S = -\frac{g_s e}{2m}\hat{S}$, where $g_s\approx 2$ is the electron $g$-factor for spin \cite{PerturbationTheory}. The total magnetic moment of an atom is the vector sum of all its orbital and spin moments. (Nuclei also have magnetic moments, but these are much smaller and often neglected in atomic magnetism.)
        Perturbation theory deals with how a system’s properties change under a small additional interaction (a perturbation). In non-degenerate perturbation theory (no two states share the same unperturbed energy), the first-order energy correction to a state is given by the expectation value of the perturbing Hamiltonian in that state. For an atomic state with no degeneracy, one can directly compute shifts due to a perturbation like a magnetic field. For example, if we apply a weak static magnetic field $\mathbf{B}$, the perturbation Hamiltonian is $H' = -\boldsymbol{\mu}\cdot \mathbf{B}$ \cite{ZeemanW}. In a non-degenerate state $|\psi^{(0)}\rangle$, first-order perturbation theory gives an energy shift $\Delta E^{(1)} = \langle \psi^{(0)}| -\boldsymbol{\mu}\cdot \mathbf{B} |\psi^{(0)}\rangle$. This is essentially the Zeeman effect: the atomic energy levels shift in proportion to the component of the magnetic moment along $\mathbf{B}$. In fact, for an electron orbital with magnetic quantum number $m_l$, one finds $\Delta E = \mu_B m_l B$ to first order \cite{Zeeman}. Including spin, the shift generalizes to $\Delta E = \mu_B g_J m_J B$, where $m_J$ is the total angular momentum projection and $g_J$ is the Landé g-factor accounting for both spin and orbital contributions \cite{ZeemanW}. Thus, non-degenerate perturbation theory predicts a linear splitting of energy levels according to their magnetic moments when a field is applied. If the unperturbed state has a non-zero magnetic moment, it will shift in energy at first order; if it has no net moment (e.g. a singlet state with $L=0$,  $S=0$), then the first-order shift can be zero and higher-order effects must be considered. Beyond energies, the wavefunction of the state is also perturbed. The first-order correction $|\psi^{(1)}\rangle$ is a sum over other (unperturbed) states, which means the perturbed state acquires a small mixture of other configurations. This can slightly alter expectation values of observables like the magnetic moment. However, for non-degenerate cases, these changes appear at least to second order in the perturbation. In summary, with non-degenerate perturbation theory we can reliably compute how a small magnetic perturbation (such as a weak field or a subtle spin interaction) shifts atomic energy levels and slightly modifies the state’s magnetic moment. This approach is valid as long as the perturbation does not connect to any degenerate levels at the same energy \cite{PerturbationTheory}.
        
    \subsection{Degenerate Perturbation Theory and Level Splitting}
        When the unperturbed atomic system has degenerate energy levels (multiple distinct states with the same energy), ordinary perturbation theory needs refinement. In degenerate perturbation theory, one must diagonalize the perturbation within the subspace of those degenerate states \cite{PerturbationTheory}. Physically, a perturbation will typically lift the degeneracy by splitting the energy levels and creating new linear combinations as eigenstates. The magnetic moment of the atom can be significantly affected by this process, since the perturbation chooses preferred orientations or couplings in the degenerate manifold. A classic example is the spin-orbit coupling inside atoms. Without spin-orbit interaction, states of the same principal quantum number $n$ with different $L$ (orbital) and $S$ (spin) orientations can be degenerate. The spin-orbit interaction $H_{S-O} = f(r), \mathbf{L}\cdot \mathbf{S}$ introduces a coupling between the electron’s orbital and spin magnetic moments \cite{PerturbationTheory}. This term mixes the degenerate $|L, m_L; S, m_S\rangle$ states. By switching to a basis of total angular momentum $J=L+S$ (i.e. coupling $L$ and $S$), one finds that $H_{S-O}$ becomes diagonal – this is effectively applying degenerate perturbation theory to find the “good” quantum states. The result is that each originally degenerate level splits into two (for $S=\tfrac{1}{2}$): one state with total $J = L+1/2$ and another with $J = L-1/2$. These have different energies given by the expectation of $\mathbf{L}\cdot \mathbf{S}$ in those states. In physical terms, the energy depends on the relative orientation of the orbital and spin magnetic moments (aligned vs. anti-aligned). One of the $J$ levels is shifted down in energy and the other up, removing the degeneracy and producing the fine structure of atomic spectra. The size of the splitting can be computed by first-order perturbation theory within the degenerate manifold, using the matrix elements of $\mathbf{L}\cdot \mathbf{S}$ \cite{PerturbationTheory}. Importantly, the new eigenstates (with definite $J$) have well-defined total magnetic moments and $g$-factors, different from the uncoupled ones, because the spin and orbital contributions are now entangled. Another example of degenerate perturbations is the Zeeman effect on a level with magnetic sublevels. In zero field, an atomic state with total angular momentum $J$ has $(2J+1)$ degenerate substates labeled by $m_J$. Turning on a magnetic field $B$ along (say) the $z$-axis adds $H_Z = -\mu\cdot B = \frac{eB}{2m}(L_z + 2S_z)$ \cite{ZeemanUT} as a perturbation. These $m_J$ states were degenerate, but $H_Z$ is diagonal in the $|J, m_J\rangle$ basis and directly splits them in energy. We can treat this by (degenerate) perturbation theory: the first-order shifts are $\Delta E = \langle J,m_J|H_Z|J,m_J\rangle$. This yields $\Delta E = \mu_B g_J B m_J$ as mentioned, so the $(2J+1)$ degenerate levels fan out into evenly spaced levels proportional to $m_J$ \cite{Zeeman}. Here the perturbation was already diagonal in the obvious degenerate basis, so the solution is straightforward (each $m$-state shifts independently without mixing). In more complicated cases (e.g. if $B$ is not aligned with a principal axis, or if there are accidental degeneracies between different $J$ levels), one would set up the Hamiltonian matrix in the degenerate subspace and diagonalize it to find the new eigenstates and energies. In all cases, degenerate perturbation theory ensures that the lifting of degeneracy is handled correctly, giving the new magnetic moment values for each split level.

    \subsection{First-Order vs. Second-Order Perturbative Effects on Magnetism}
        In perturbation theory, effects can appear at different orders in the perturbation strength (often denoted by a power of a small parameter). First-order effects are typically linear in the perturbation, while second-order effects are proportional to the perturbation squared (or involve two perturbative interactions). Both can influence atomic magnetic properties in different ways:
        \begin{enumerate}
            \item First-Order Perturbations: These give the dominant, linear response. If an atomic state has a permanent magnetic moment (e.g. due to unpaired spin or orbital angular momentum), a first-order perturbation will usually produce a direct change in energy or orientation. For instance, the Zeeman splitting discussed above is a first-order effect: energy levels shift linearly with $B$ according to $\Delta E^{(1)} \propto B$ \cite{Zeeman}. Likewise, the spin-orbit coupling term $H_{S-O} = f(r),\mathbf{L}\cdot\mathbf{S}$ produces a first-order fine-structure splitting of atomic levels proportional to the strength of that interaction \cite{PerturbationTheory}. First-order perturbations can also change the wavefunction at order $\lambda^1$, which in turn can lead to first-order changes in observable expectation values if the unperturbed expectation was non-zero. In general, first-order effects account for things like the immediate alignment of magnetic dipoles with a field (paramagnetism) and the primary shifts in spectra.
            Second-Order Perturbations: These are typically smaller corrections, but they can become important in cases where first-order effects cancel or are forbidden. Second-order energy corrections involve virtual transitions to intermediate states and are proportional to (perturbation strength)$^2$. In magnetism, second-order perturbations can induce a magnetic moment even when the atom has none initially, or they can modify the first-order results. A good example is the diamagnetic response of closed-shell atoms: if an atom has $J=0$ (total angular momentum zero) in its ground state, the first-order Zeeman shift is zero (since $m_J=0$ gives $\Delta E^{(1)}=0$). However, a magnetic field will still induce currents in the electron cloud. This appears as a second-order effect: the field mixes the ground state with excited states, leading to a small induced magnetic moment opposite to the field (Lenz’s law) and an energy shift proportional to $B^2$ (the quadratic Zeeman effect). For exam

ple, in the hydrogen 1s state (which has no orbital angular momentum), the leading Zeeman effect comes from a second-order term resulting in $\Delta E \propto +B^2$ (a small diamagnetic shift)

 \cite{QZeeman}. Second-order perturbation theory also explains Van Vleck paramagnetism, a subtle form of magnetism in certain ions: there, a $J=0$ ground state can have a weak, temperature-independent paramagnetism due to excited-state mixing in second order \cite{Magnetism}. Generally, second-order perturbations refine the magnetic properties by accounting for induced moments and polarization effects, and they become crucial when first-order terms vanish or when higher precision is needed. They also give corrections like the quadratic Zeeman splitting (important at high fields or in precision spectroscopy) and small shifts in $g$-factors due to state mixing.
        \end{enumerate}
 

\section{


# Suppression of Bremsstrahlung in Relativistic p–¹¹B Fusion Plasmas

## Background: Bremsstrahlung Losses in Proton–Boron Fusion

Proton–boron-11 (p–¹¹B) fusion is an attractive aneutronic reaction, but its feasibility is limited by severe radiative energy losses. High-energy electrons in the hot plasma emit copious bremsstrahlung X-rays when deflected by ionic Coulomb fields, often outpacing the fusion power output. In a relativistic plasma (with electron temperatures of order hundreds of keV), bremsstrahlung emission is enhanced by the long high-energy tail of the electron distribution. Unlike lower-temperature plasmas, **electron–electron** collisions also contribute significantly to bremsstrahlung at relativistic energies, adding to the usual **electron–ion** bremsstrahlung losses. These losses pose a fundamental challenge: a classical Maxwellian p–¹¹B plasma would radiate a large fraction of its energy away, preventing net power gain. Reducing or suppressing such radiation, especially bremsstrahlung, is therefore critical for p–¹¹B fusion viability.

## Relativistic Bremsstrahlung Theory in Strong Fields

At relativistic energies, the bremsstrahlung differential cross-section (Bethe–Heitler formula) grows with electron energy, meaning **superthermal electrons dominate X-ray emission**. This has two important implications for suppression strategies. First, *shaping the electron energy distribution* to minimize the high-energy tail can disproportionately cut radiative losses. Second, intense electromagnetic fields can alter electron trajectories and quantum states, modifying emission probabilities. For example, in ultra-strong magnetic fields, electrons occupy discrete Landau energy levels (quantized transverse motion), which changes collision dynamics. If the field is extremely high (above \~5×10^9 G), ion–electron collisions can no longer easily excite electrons to higher Landau levels. In effect, **Landau quantization raises the threshold for energy transfer to electrons**, keeping electron temperatures lower relative to ions. This can dramatically reduce bremsstrahlung cooling in principle, as theorized for magnetized proton–boron plasmas where fusion power could far exceed X-ray losses under multi-gigagauss fields. Such quantum magnetic-field effects (Landau quantization and related QED phenomena) form an extreme end of bremsstrahlung suppression – likely relevant in astrophysical or dense plasma focus scenarios, but requiring magnetic fields beyond conventional lab scales.

Strong magnetic fields also introduce classical suppression effects. In relativistically intense laser–plasma interactions, quasi-static magnetic fields on the order of tens of megatesla are predicted to reduce bremsstrahlung emission probability by altering electron orbits. This *magnetic suppression mechanism* occurs because a large magnetic field bends electron trajectories, effectively limiting the duration and straight-line distance over which an accelerating electron can emit a bremsstrahlung photon. **Habibi et al. (2023)** implemented this effect in Particle-in-Cell (PIC) simulations using the *EPOCH* code, incorporating magnetic and electric field influences into the bremsstrahlung emission module. Their study – the first of its kind – showed that a strong field suppresses especially the low-energy X-ray emissions and also alters the dynamics of radiating electrons. This suggests that in magnetized relativistic plasmas, bremsstrahlung can be partially mitigated by field-induced trajectory perturbations, complementing any quantum Landau-level effects.

## Tailoring Electron Distributions (Energy Cutoff Method)

Because high-energy electrons disproportionately drive bremsstrahlung, one intuitive approach is to **remove or redistribute the electron tail** to lower energies. This idea can be realized by introducing an *energy cutoff* in the electron distribution. Instead of a pure Maxwell–Jüttner distribution, the plasma electrons follow a distribution where beyond some cutoff energy \$E\_c\$ the population is sharply reduced. Recent simulations and analytic models have explored how such cutoffs suppress radiation. **Munirov and Fisch (2023)** showed that *redistributing superthermal electrons into lower energies can significantly reduce bremsstrahlung losses in a relativistic plasma*. Importantly, this suppression effect is a uniquely relativistic benefit – in nonrelativistic plasmas removing the electron tail has much less impact on total radiation. Munirov and Fisch calculated bremsstrahlung power for a range of cutoff energies and electron temperatures, keeping total electron density fixed (i.e. electrons are not removed from the system but slowed down). They found that even a modest cutoff can yield measurable reductions. For example, at an electron temperature \$T\_e \sim 100\$ keV (typical of p–¹¹B reactors), imposing a cutoff around \$E\_c \approx 70\$–100 keV can **lower electron–ion bremsstrahlung power by \~10% and electron–electron bremsstrahlung by \~20%** relative to a Maxwellian plasma. Higher cutoffs (>2–3 times the thermal energy) see diminishing returns, as most electrons remain below the cutoff.

There are several physical mechanisms to achieve such an effective energy cutoff. In *open magnetic field-line devices* (like mirror machines or certain inertial-electrostatic confinement setups), high-energy electrons can be allowed to escape along field lines, while colder electrons and ions are confined by electrostatic potentials. The escaping hot electrons carry kinetic energy out, effectively truncating the distribution inside the core plasma. As one example, a **magnetic mirror** or ambipolar electrostatic well can be tuned so that electrons above a certain energy have enough parallel velocity to overcome the confining potential and stream out, whereas lower-energy electrons remain trapped. This *selective loss* mechanism continuously removes the nascent fast electrons as they are produced by collisions, preventing a high-energy tail from accumulating. The lost electrons’ charge is compensated by ion outflow or by replenishing with slower electrons to maintain quasineutrality. Overall, the plasma sustains a skewed distribution with far fewer relativistic electrons than a Maxwellian, thereby emitting less bremsstrahlung.

Another approach to distribution control is through **waves and instabilities** that damp high-energy electrons. *Resonant cyclotron emission* can act as a radiative sink for electrons above a threshold energy, as noted in strong-field plasma calculations. In a 1999 study, Riffert *et al.* found that in a magnetized proton–electron plasma, electrons above the cyclotron resonance energy rapidly lose energy to **intense cyclotron (synchrotron) radiation**, depleting the high-energy tail and leaving a narrow spectral line at the cyclotron frequency. This indicates that under certain conditions, **synchrotron cooling can outpace bremsstrahlung for the fastest electrons**, effectively capping their energy and reducing continuum X-ray emission. Similarly, modern PIC simulations of laser-driven plasmas have observed that including radiation reaction forces (synchrotron losses) leads to a natural high-energy cutoff in the electron spectrum. While synchrotron emission is itself an energy loss channel, it tends to be more peaked (narrow-band) and may be easier to capture or mitigate than broadband bremsstrahlung. This interplay of radiation processes offers multiple avenues to engineer the electron distribution for minimal bremsstrahlung: either by **external fields that allow fast electron escape** or by **internal plasma processes that selectively drain energy from fast electrons**.

## Resonant Magnetic Perturbations and Stochastic Confinement

Magnetic perturbation techniques exploit three-dimensional fields to alter confinement, and they can be a powerful tool for controlling high-energy particles. In toroidal fusion devices (tokamaks and stellarators), applying *resonant magnetic perturbations (RMPs)* has been well-studied as a method to induce magnetic islands and stochastic regions in the plasma edge. These perturbations, typically small static fields introduced by external coils, resonate with the plasma’s field geometry and partially break the confining magnetic surfaces. The resulting **magnetic islands and open field lines** can allow particles (and heat) to flow out of the plasma more readily. While RMPs are traditionally used for edge-localized mode (ELM) control, recent work indicates they can also *enhance the loss of runaway electrons* and other suprathermal particles – directly relevant to bremsstrahlung suppression since runaways generate intense X-ray emission.

Simulation studies have examined how different RMP spectra affect fast electron confinement. **Cornille *et al.* (2022)** performed MHD and test-particle simulations of runaway electron behavior in the MST tokamak with applied RMPs. Using the NIMROD code (for nonlinear MHD) coupled with tracer electrons (using the KORC orbit-following code), they found that an RMP with a certain mode structure can **chaoticize the magnetic field in the plasma edge and substantially deconfine high-energy electrons**. In their results, an \$m=3\$ RMP (with three island chains around the torus) induced a stochastic outer region that led to \~50% of runaway electrons being lost from the plasma, compared to only \~10–15% loss without RMP or with a non-resonant \$m=1\$ perturbation. The RMP effectively opened up escape paths for the energetic electrons, preventing a fully developed runaway beam. These findings are consistent with experimental observations on devices like J-TEXT and TEXTOR, where applying resonant \$n=1\$ or \$n=2\$ perturbation fields during disruptions measurably increased runaway electron losses (sometimes achieving full suppression of a runaway current). In the fusion plasma context, **RMPs introduce a controllable level of magnetic chaos**: by tuning the perturbation amplitude and spectrum, one can strike a balance between purging high-energy electrons and still confining the bulk plasma. This concept directly parallels the energy-cutoff approach – instead of relying on ambipolar electrostatic fields as in mirror machines, a tokamak could use magnetic topological tweaks to let the most energetic electrons escape (or hit a limiter/divertor) before they radiate.

State-of-the-art simulation tools support these studies of 3D field effects. The *JOREK* non-linear MHD code, for instance, has been used to model plasma response to RMPs in tokamaks, capturing the formation of magnetic islands and the resultant changes in transport. JOREK and similar codes (e.g. NIMROD, M3D-C1) can include drift-kinetic or fluid models for runaway electrons, allowing integrated simulations of how applied perturbations dissipate fast electron populations. Meanwhile, *HINT/HINT2* (an MHD equilibrium solver) and analytic field-line tracing codes like *FLARE* provide ways to visualize and design magnetic configurations with open field lines. **FLARE** in particular can compute Poincaré maps and connection lengths in non-axisymmetric fields, helping researchers identify how an RMP or helical field will intersect material walls or direct particles to a divertor target. Using these tools, one can simulate a scenario where, for example, **slow electrons remain confined on closed flux surfaces while fast electrons wander into stochastic regions and are removed** – effectively creating the desired energy-selective confinement. Such *magnetic perturbation-based techniques* are a promising avenue for bremsstrahlung suppression in advanced p–¹¹B reactors, especially in designs that start from a toroidal confinement concept and need to shed excess electron energy.

## Quantum Magnetic Effects: Landau Quantization and Radiative Decoupling

For completeness, it is worth emphasizing the role of extreme magnetic fields and quantum effects on radiative processes. In magnetized plasmas approaching the quantum limit (where the cyclotron energy \$\hbar \omega\_c\$ is comparable to electron thermal energy), electrons occupy quantized Landau levels for their gyration. This **Landau quantization** alters plasma microphysics in ways that can reduce radiation. As mentioned earlier, one effect is suppressed collisional heating of electrons: if ions cannot easily excite electrons to higher Landau levels, the electrons remain cooler, and bremsstrahlung emission (which scales strongly with electron energy) is curtailed. Another effect is that an electron in a Landau-quantized state has restricted degrees of freedom for radiative transitions. For instance, direct bremsstrahlung emission might be reduced because the electron’s transverse momentum changes are quantized – in essence, the phase space for emitting a photon is constrained by selection rules. Although a full quantum electrodynamics treatment is complex, studies in neutron star and magnetar physics suggest that **bremsstrahlung and even pair-production rates are significantly modified in ultra-strong fields**. In fusion-oriented research, Lerner and colleagues argued that fields on the order of several gigagauss could *decouple electron and ion temperatures* and thereby mitigate X-ray losses in p–¹¹B plasmas. Achieving such fields in the lab is challenging; however, pulsed power devices like the dense plasma focus (DPF) have reported multi-megagauss (several \$10^8\$ G) fields in compact plasmoids, and pulsed magnets or laser-driven capacitor coils have reached the 10–100 MG range (1–10 ×10^7 G) in short bursts. At those field strengths, some level of Landau quantization begins to appear (the cyclotron energy for electrons becomes tens of keV). **If future high-field schemes can push towards the 1 GG (10^9 G) scale**, we may enter a regime where bremsstrahlung suppression by Landau quantization is noticeable – complementing classical techniques by fundamentally changing how electrons radiate.

## Simulation Frameworks for Radiation Suppression Studies

Research into these suppression techniques leans heavily on simulations, as many scenarios are beyond current experimental reach. Various computational models are employed, each with strengths for capturing different physics:

* **Particle-in-Cell (PIC) simulations:** PIC codes like *EPOCH*, *OSIRIS*, and *Smilei* directly track particles and electromagnetic fields. They are well-suited to model kinetic effects such as anisotropic tail formation, wave–particle interactions, and radiation reaction in laser-driven or magnetic-fusion plasmas. For example, Habibi *et al.* added a bremsstrahlung module to EPOCH to study radiation with and without magnetic suppression. PIC simulations can include Monte Carlo bremsstrahlung emission and even QED processes (photon emission, pair creation) in extreme fields, providing a first-principles view of how radiation is produced and can be mitigated by field effects.

* **Magnetohydrodynamic (MHD) and hybrid codes:** Nonlinear MHD codes such as *JOREK*, *NIMROD*, and *M3D-C1* treat the bulk plasma as a fluid but can be coupled with kinetic markers or reduced models for fast particles. These codes are crucial for studying global confinement changes due to RMPs or other perturbations. As shown in the MST tokamak simulation, an MHD code computed the perturbed magnetic topology while a particle tracker computed electron losses. JOREK has recently been extended with a runaway electron fluid model to simulate how an applied 3D field or disruption affects runaway generation and loss, which ties directly into bremsstrahlung outcomes. Such **hybrid simulations** (fluid plasma + kinetic superparticles for runaways or alphas) allow exploration of scenarios like ITER where avoiding runaway-driven bremsstrahlung during disruptions is critical.

* **Field line tracing and equilibrium codes:** As mentioned, codes like *HINT2* (for 3D equilibrium) and *FLARE* (for field-line reconstruction) provide the magnetic geometry underpinning transport simulations. These tools don’t simulate plasma dynamics directly but are used to set up initial conditions or analyze the results – for example, computing the connection length of field lines in an RMP perturbed plasma to see how deeply the perturbation penetrates. Long connection lengths indicate confined regions, whereas finite lengths indicate where particles will hit walls/divertors. By correlating these with test-particle simulations, researchers can identify which perturbation harmonics preferentially eject high-energy electrons while minimally perturbing thermal plasma confinement.

* **High-field and QED models:** In regimes approaching QED limits (intense laser–plasma or astrophysical conditions), specialized codes incorporate quantum effects. Some PIC codes (e.g. EPOCH, PICLS) have modules for Landau quantization and radiation reaction. Additionally, **Monte Carlo codes** like PENELOPE and Geant4 can be adapted to simulate bremsstrahlung with suppression factors (e.g. the Landau-Pomeranchuk-Migdal effect in materials, or magnetic suppression as a function of field strength). These are generally used to post-process or validate analytical predictions for radiation yields in strong fields.

## Toward a Thesis: Integrating Suppression Strategies for p–¹¹B Fusion

In summary, a range of simulation studies indicate that bremsstrahlung in relativistic p–¹¹B fusion plasmas is not an immutable “show-stopper” but can be **actively suppressed through intelligent plasma control**. Magnetic perturbation-based techniques emerge as especially promising. By leveraging RMPs or open-field configurations, one can enforce an **effective energy cutoff** in the electron population, significantly reducing radiative power loss. The lost energy is carried by escaping electrons – which, unlike X-rays, are charged and can be captured or converted (for instance, via direct electricity converters or beam dumps). This traded loss channel is more manageable than isotropic bremsstrahlung, which would otherwise require surrounding the reactor with heavy radiation shielding and energy recovery systems.

Furthermore, extreme high-field approaches using modern pulsed-power technology hint that **quantum suppressive effects** (Landau level physics) could reinforce these classical strategies. Even a 10–20% reduction in bremsstrahlung, as predicted by moderate cutoff models, might tip the balance for p–¹¹B reactors toward net positive energy, especially when combined with other innovations like alpha channeling and advanced direct conversion. Going forward, **multi-physics simulation platforms** (coupling PIC, MHD, and kinetic theory) will be essential to design and optimize reactors that exploit these suppression mechanisms. The literature reviewed – from PIC experiments of magnetic suppression to MHD simulations of RMP-driven losses – provides a strong theoretical grounding. This will feed into a Master’s thesis proposal focused on *novel bremsstrahlung suppression in aneutronic fusion*, where the candidate will propose to use and extend codes like EPOCH, JOREK, and FLARE to evaluate integrated scenarios. By building on these cutting-edge studies, the thesis can chart a path toward a plasma regime in which **radiative losses are tamed**, bringing p–¹¹B fusion a step closer to practicality.

    
\section{Research design and methods}
\section{Potential Implications}
\printbibliography    
\end{document}
